\documentclass[11pt,a4paper]{article}

\usepackage[T1]{fontenc}
\usepackage{lmodern}
\usepackage[margin=28mm]{geometry}
\usepackage{amsmath,amssymb,amsthm}
\usepackage{booktabs}
\usepackage{microtype}
\usepackage[hidelinks]{hyperref}

\newtheorem{theorem}{Theorem}
\newtheorem{lemma}[theorem]{Lemma}
\newtheorem{corollary}[theorem]{Corollary}
\newtheorem{proposition}[theorem]{Proposition}
\theoremstyle{definition}
\newtheorem{definition}[theorem]{Definition}
\theoremstyle{remark}
\newtheorem{remark}[theorem]{Remark}
\newcommand{\Ostar}{O^{*}}
\newcommand{\E}{\mathbb{E}}
\newcommand{\sat}{\mathrm{sat}}

\clubpenalty=10000
\widowpenalty=10000
\displaywidowpenalty=10000
\emergencystretch=2em
\hypersetup{pdftitle={Max 2-SAT on Unaligned Formulas}}

\title{Max 2-SAT Below $2^n$ on Unaligned Formulas:\\ Theorem and Proof Sketch}
\author{}
\date{8 October 2026\\[3pt]\small Working notes}

\begin{document}
\maketitle

\begin{abstract}
We define an explicit class of 2-CNF formulas by three conditions on the
quadratic polynomial that counts satisfied clauses. Every formula in the class
is solved exactly by a deterministic polynomial-space algorithm in time
$\Ostar(c^n)$, with a constant $c<2$ that depends only on a parameter of the
class and not on the number of clauses. The class contains almost all formulas
of density at least $C\log n$ and excludes the aligned formulas that defeat
single-flip pruning. The algorithm is the search of~\cite{N7}: pruning clauses
and the sum check, along a fixed order. The proof is the argument of~\cite{N8}
for a fixed formula: random signs are replaced by the conditions of the class,
and the only randomness left is the random partial assignment of Knuth's
identity. This note states the theorem and sketches the proof; the details
have not been written out.
\end{abstract}

\section{The class and the theorem}\label{sec:theorem}

Let $F$ be a 2-CNF formula over $x_1,\dots,x_n$. Its clauses have one literal
or two literals on distinct variables; several clauses on the same pair of
variables are allowed, but no binary clause occurs twice. A literal on $x_v$
has sign $s=+1$ if it is $x_v$ and $s=-1$ if it is $\neg x_v$. For an
assignment write $a_v=+1$ if $x_v$ is true and $a_v=-1$ otherwise.

\begin{definition}[Coefficients]\label{def:coef}
\begin{enumerate}
 \item The \emph{parity matrix} of $F$ is the symmetric $n\times n$ integer
 matrix $J=(J_{uv})_{u,v=1}^n$ with zero diagonal, $J_{vv}=0$, and, for
 $u\ne v$,
 \[
 J_{uv}=\sum_{\text{binary clauses on }\{x_u,x_v\}}s_us_v ,
 \]
 where $s_u,s_v$ are the signs of the two literals of the clause. So
 $J_{uv}=0$ if no clause contains both $x_u$ and $x_v$.
 \item The \emph{imbalance vector} of $F$ is $h=(h_1,\dots,h_n)\in\mathbb
 Z^n$, where $h_v$ is the sum of the signs of $x_v$ over its binary clauses
 plus twice the sum of its signs over its unit clauses.
 \item $L$ is the strictly lower triangular part of $J$: the $n\times n$
 matrix with $L_{uv}=J_{uv}$ for $u>v$ and $L_{uv}=0$ for $u\le v$. It
 depends on the numbering of the variables.
 \item $d=\frac1n\sum_{u,v}J_{uv}^2$, the mean of the squared row norms of
 $J$.
\end{enumerate}
\end{definition}

A binary clause $(\ell_u\lor\ell_v)$ is satisfied with indicator
$\frac34+\frac14s_ua_u+\frac14s_va_v-\frac14s_us_va_ua_v$, and a unit clause
$(\ell_v)$ with indicator $\frac12+\frac12s_va_v$. Hence the number of
satisfied clauses is
\[
 \sat_F(a)=c_0+\frac14\sum_vh_va_v-\frac14\sum_{u<v}J_{uv}a_ua_v ,
\]
where $c_0$ is $\frac34$ times the number of binary clauses plus $\frac12$
times the number of unit clauses. So $J$ and $h$ are the coefficients of the score, and every simplification that
changes the score by a constant leaves them unchanged. For example,
$(x\lor y)$ and $(x\lor\neg y)$ together score $1+[x]$: they give $J_{xy}=0$
and add $2$ to $h_x$, exactly as the unit clause $(x)$ does. A complementary
pair $(x\lor y),(\neg x\lor\neg y)$ gives $J_{xy}=2$ and adds nothing to $h$.
If every pair carries at most one clause, $J$ is the signed adjacency matrix
with entries $\pm1$ and $d=2m/n$ is the average degree.

\begin{definition}[Unaligned formulas]\label{def:unaligned}
For $K\ge1$, a formula is \emph{$K$-unaligned} if
\[
 \|J\|\le K\sqrt d,\qquad \|L\|\le K\sqrt d,\qquad \sum_vh_v^2\le K\,nd,
\]
where $\|\cdot\|$ is the spectral norm.
\end{definition}

\begin{theorem}\label{thm:main}
For every $K\ge1$ there is $\varepsilon_K>0$, for instance
$\varepsilon_K=10^{-5}K^{-4}$, such that Max 2-SAT on $K$-unaligned formulas
can be solved by a deterministic algorithm in time
$\Ostar(2^{(1-\varepsilon_K)n})$ using polynomial space, for every number of
clauses.
\end{theorem}

The algorithm is the same for all $K$; it does not need to know $K$ or to test
membership in the class.

\begin{remark}[About the class]\label{rem:class}
The class has the following properties.
\begin{enumerate}
 \item Membership can be checked in polynomial time.
 \item Random formulas lie in the class with small $K$
 (Table~\ref{tab:params}). If each pair carries a clause with probability
 $d/(n-1)$, with uniform signs, and $d\ge C\log n$, then the formula is
 $3$-unaligned with probability $1-o(1)$. This follows from norm bounds for
 random matrices with independent entries~\cite{BvH16} and from Chernoff
 bounds for $\sum_vh_v^2$.
 \item The aligned formulas of~\cite{N3,N4} lie far outside. In the circulant
 formulas $C_{n,c}$ all entries of $J$ equal $1$, and in Max-Cut on $K_m$ all
 equal $2$, so $\|J\|$ is of order $d$ and $K$ would have to be of order
 $\sqrt d$.
 \item The numbering of the variables enters only through $L$; the algorithm
 branches in this order.
 \item Sparse random formulas with $d$ below about $\log n$ are not
 $K$-unaligned for a fixed $K$, because a few variables of high degree make
 $\|J\|$ large. For them only the average-case result of~\cite{N8} is
 available.
 \item The definition is spectral; the algorithm is not.
\end{enumerate}
\end{remark}

\begin{table}[ht]
\centering
\begin{tabular}{lcccc}
\toprule
 & $d$ & $\|J\|/\sqrt d$ & $\|L\|/\sqrt d$ & $\sum_vh_v^2/(nd)$\\
\midrule
random, one clause per pair & 30 & 2.02 & 1.67 & 1.02\\
 & 100 & 2.03 & 1.66 & 1.01\\
 & 999 & 2.00 & 1.64 & 0.94\\
random, pairs and unit clauses & 189 & 2.00 & 1.66 & 0.38\\
circulant $C_{401,100}$ & 400 & 20.0 & 12.8 & 0.00\\
Max-Cut on $K_{400}$ & 1596 & 20.0 & 12.7 & 0.00\\
\bottomrule
\end{tabular}
\caption{The three parameters of Definition~\ref{def:unaligned}; random
formulas have $n=1000$. In the fourth row, $30\%$ of the interacting pairs are
complementary pairs of random type, and each variable has up to two random
unit clauses.}\label{tab:params}
\end{table}

\section{The algorithm}\label{sec:algorithm}

\paragraph{Preprocessing.} Exact pair reductions~\cite{N3} bring $F$ into
canonical form: unit clauses, kept with multiplicity, and on each pair of
variables nothing, one binary clause, or one complementary pair. They change
the score by a constant, so $J$, $h$, $L$ and $d$ do not change.

\paragraph{The search~\cite{N7}.} Branch on $x_1,\dots,x_n$ in this order. At
a node where $x_1,\dots,x_{t-1}$ are assigned, let $\Delta_t$ be the effective
gap of $x_t$: the sum of its signs over its remaining binary clauses plus twice
the sum over its unit clauses. Let $D_t$ be the number of its remaining binary
clauses, $\ell$ its preferred literal ($x_t$ if $\Delta_t\ge0$, else $\neg
x_t$), and $p=(D_t-|\Delta_t|)/2$. The child in which $\ell$ is false receives
the pruning clause $\langle N_t\le p-1\rangle$. Here $N_t$ counts the true
literals among $\neg m$ for each binary clause $(\ell\lor m)$ and $m$ for each
binary clause $(\neg\ell\lor m)$, with multiplicity; a complementary pair
contributes the same literal twice. The bad value of a variable in a pruning
clause is the value that makes its literal true. A node is cut when a pruning
clause has more true literals than it allows, or when the \emph{sum check}
fires:
\[
 \sum_{w\text{ free}}\min\{m^0_w,m^1_w\}>\sum_ck_c ,
\]
over the live pruning clauses that are not vacuous. Here $k_c$ is the remaining
allowance of $c$, and $m^b_w$ counts, with multiplicity, the clauses in which
$b$ is the bad value of $w$.

\paragraph{Soundness and cost.} Each branching adds a clause to one child
only, and the flip identity holds for canonical formulas with multiplicities.
So the path-order argument of~\cite[Proposition~20]{N3} shows that an optimal
assignment survives, and the sum check is implied by the clauses~\cite{N4}.
Let $\pi_t$ be the probability that a uniformly random $\sigma\in\{0,1\}^n$
survives the checks at depths $1,\dots,t$. By Knuth's
identity~\cite{Knu75,N7}, the search has at most $1+2n\max_t2^t\pi_t$ nodes. It
uses polynomial space and polynomial time per node.

\section{Proof sketch}\label{sec:proof}

Fix a $K$-unaligned formula; all probabilities are over the uniform $\sigma$.
Put $a=2\sigma-1$ and $T=(1-\delta)n$, with $\delta=\delta(K)$ chosen in
Step~6.

\paragraph{Step 1: one depth.} $\max_t2^t\pi_t\le\max\{2^T,\,2^n\pi_T\}$,
because $\pi_t$ is non-increasing. It suffices to show $\pi_T\le6e^{-cn}$.

\paragraph{Step 2: the race.} Run $\sigma$ through the search without cutting
and keep every pruning clause created up to depth $T$. Put
$Z_T=S_T+Q_T-U_T$, with
\[
 S_T=\sum_c\Bigl(\frac{|\Delta_c|}2+1\Bigr),\quad
 Q_T=\sum_c\Bigl(j_c-\frac{e_c}2\Bigr),\quad
 U_T=\frac12\sum_{w>T}|m^1_w-m^0_w| .
\]
Here $\Delta_c$ is the gap at the creation of $c$, $e_c$ is the number of its
assigned support variables, and $j_c$ is the number of those that took their
bad value, both with multiplicity. By the race identity and the remark on
vacuous clauses~\cite[Lemma~2, Remark~3]{N7}, whose algebra is unchanged by
multiplicities, a $\sigma$ that survives depth $T$ has $Z_T\le0$.

\paragraph{Step 3: the race in terms of $J$ and $h$.}
\begin{lemma}[Representation]\label{lem:rep}
Let $C$ be the set of $s\le T$ at which $\sigma_s$ is the non-preferred value,
and $y_s=a_s$ for $s\in C$, $y_s=0$ otherwise. Then
\begin{enumerate}
 \item[(a)] $\Delta_t=h_t-\sum_{u<t}J_{tu}a_u$, and $t\in C$ if and only if
 $a_t=-\operatorname{sign}\Delta_t$ (with $\operatorname{sign}0=+1$);
 \item[(b)] the balance $m^1_w-m^0_w$ of a variable $x_w$, taken when $x_w$ is
 assigned (or at depth $T$ if $w>T$), equals $\sum_{s<\min\{w,T+1\}}J_{sw}y_s$;
 \item[(c)] $S_T=\sum_{t\in C}(|\Delta_t|/2+1)$,
 $Q_T=\frac12\sum_{t\le T}a_t(Ly)_t$ and
 $U_T=\frac12\sum_{w>T}|\sum_{s\le T}J_{sw}y_s|$.
\end{enumerate}
\end{lemma}

\noindent\emph{Why.} After $a_1,\dots,a_{t-1}$ are fixed, the coefficient of
$a_t$ in the score polynomial is $\frac14(h_t-\sum_{u<t}J_{tu}a_u)$. The
remaining formula is canonical, and there the same coefficient is
$\frac14\Delta_t$; this gives~(a). For~(b), let $x_s$ create a clause and $x_w$
be in its support. A case distinction shows that the bad value of $x_w$,
written as $\pm1$ and taken with its multiplicity, is $J_{sw}a_s$. For a single
clause this is $\pm1$. For a complementary pair both clauses give the same bad
value, $\sigma_s$ for the pair $(x\lor y),(\neg x\lor\neg y)$ and
$1-\sigma_s$ for the other type, which is $J_{sw}a_s$ with $J_{sw}=\pm2$.
Part~(c) follows: $Q$ changes at step $t$ by $\frac12a_t(m^1_{x_t}-m^0_{x_t})$.
All parts were checked numerically, including complementary pairs and unit
clauses (Section~\ref{sec:checks}).

\paragraph{Step 4: uncertainty, for every $\sigma$.} By (c) and
Cauchy--Schwarz, using $\|y\|^2\le T$,
\[
 U_T\le\tfrac12\sqrt{n-T}\,\|J\|\,\|y\|\le\tfrac12\sqrt{T(n-T)}\,\|J\|
 \le\tfrac12K\sqrt{\delta}\,n\sqrt d .
\]

\paragraph{Step 5: the martingale.} By (c), $Q$ changes at step $t$ by
$\pm\frac12(Ly)_t$. The sign is $a_t$, which is independent of
$\sigma_1,\dots,\sigma_{t-1}$, while $(Ly)_t$ is determined by them. The squared
steps sum to at most $\frac14\|L\|^2\|y\|^2\le\frac14K^2dn$ for every $\sigma$.
Azuma's argument gives
$\Pr[Q_T\le-\kappa]\le\exp(-2\kappa^2/(K^2dn))$.

\paragraph{Step 6: the shortfall.} Let $G=\sum_{t\le T}|\Delta_t|$.
\begin{enumerate}
 \item[(a)] \emph{Coin flips.} Given $\sigma_1,\dots,\sigma_{t-1}$, the gap
 $\Delta_t$ is fixed and $t\in C$ with probability $\frac12$. So
 $X_t=[t\in C]\,|\Delta_t|/2\ge0$ has conditional mean $|\Delta_t|/4$ and
 conditional second moment $\Delta_t^2/8$. Moreover
 $\sum_t\Delta_t^2\le2\|h\|^2+2\|L\|^2n\le2(K+K^2)nd$ by~(a). The one-sided
 bound for nonnegative increments~\cite{Mau03,N8} gives
 $S_T\ge\frac14G-\alpha$, except with probability
 $\exp(-\alpha^2/(2V))$, where $V=(K+K^2)nd/4$.
 \item[(b)] \emph{Mean.} By (a), $\Delta_t$ is a fixed number plus the
 symmetric sum $-\sum_{u<t}J_{tu}a_u$. So by symmetry and Szarek's
 inequality~\cite{Sza76}, $\E|\Delta_t|\ge(W_t/2)^{1/2}$ with
 $W_t=\sum_{u<t}J_{tu}^2$. Every row of $J$ has
 $\sum_uJ_{tu}^2\le\|J\|^2\le K^2d$, so
 $\sum_{t\le T}W_t\ge\frac12nd-\delta n K^2d$ and $W_t\le K^2d$. Hence
 \[
 \E G\ge\sum_{t\le T}\sqrt{W_t/2}\ge\frac{\sum_tW_t}{\sqrt2\,K\sqrt d}
 \ge\frac{\frac12-\delta K^2}{\sqrt2\,K}\,n\sqrt d .
 \]
 \item[(c)] \emph{Concentration.} $G$ is a convex function of
 $a\in[-1,1]^n$, Lipschitz with constant $\sqrt T\|L\|\le K\sqrt{nd}$. By
 Talagrand's inequality~\cite{Tal95},
 $\Pr[G\le\E G-15K\sqrt{nd}-\eta]\le4\exp(-\eta^2/(16K^2nd))$. Here
 $15K\sqrt{nd}$ bounds the distance between the mean and the median.
\end{enumerate}

\paragraph{Step 7: the constants.} Take $\delta=1/(128K^4)$. Steps 4 and 6
give a margin
\[
 A=\frac{\frac12-\delta K^2}{4\sqrt2\,K}-\frac K2\sqrt\delta\ \ge\
 \frac{0.0428}K
\]
in units of $n\sqrt d$. Take $\alpha=\kappa=\eta=\frac14An\sqrt d$. Outside the
three exceptional events,
\[
 Z_T\ \ge\ \tfrac14(\E G-15K\sqrt{nd}-\eta)-\alpha-\kappa-U_T\
 \ge\ \tfrac7{16}An\sqrt d-4K\sqrt{nd}\ >\ 0
\]
for $n\ge n_0(K)$. The three events have probabilities at most
$e^{-A^2n/(16K^2)}$, $4e^{-A^2n/(256K^2)}$ and $e^{-A^2n/(8K^2)}$. So
$\pi_T\le6e^{-cn}$ with $c=A^2/(256K^2)\ge7\cdot10^{-6}K^{-4}$. By Step~1 the
search has $\Ostar(2^{(1-\varepsilon)n})$ nodes, with
$\varepsilon=\min\{\delta,c/\ln2\}\ge10^{-5}K^{-4}$. \hfill$\square$

\section{What has to be written out}\label{sec:todo}

\begin{enumerate}
 \item The race identity and the remark on vacuous clauses with
 multiplicities, and Lemma~\ref{lem:rep} in full. All are checked numerically,
 but only the version with one clause per pair is written~\cite{N7,N8}.
 \item The exact form of Talagrand's inequality used in Step~6(c), with its
 constants.
 \item The claim of Remark~\ref{rem:class}(2) on random formulas, with the
 constants of~\cite{BvH16}.
\end{enumerate}
The constants are not optimized. The measured saving of the same search on
random formulas is about $0.5n$~\cite{N7}, many orders of magnitude more than
$\varepsilon_K$. Three questions stay open: whether $\|L\|$ can be replaced by
$\|J\|$ (it is at most $O(\log n)$ times larger, and that factor would cost
too much); whether the class can be widened to sparse formulas by discounting
high-degree variables; and whether repeated clauses (weights) can be admitted.

\section{Checks}\label{sec:checks}

\texttt{explore\_explicit\_class.py}, each mode in a few seconds:
\begin{itemize}
 \item \texttt{identities}: $128{,}826$ checks on random canonical formulas
 (one clause per pair, complementary pairs, unit clauses) and on $C_{21,6}$,
 $C_{31,9}$, Max-Cut on $K_{12}$ and $K_{16}$. Covered: the gap formula and
 the bad-value formula of Lemma~\ref{lem:rep}, the race identity with
 multiplicities, the formula for $U_T$, and $Z_t\le0$ on every surviving
 prefix.
 \item \texttt{params}: Table~\ref{tab:params}.
 \item \texttt{constant}: Step~7 for $K=1,\dots,4$, giving
 $\varepsilon_K K^4\approx1.04\cdot10^{-5}$ to $1.10\cdot10^{-5}$.
\end{itemize}

\begin{thebibliography}{99}
\raggedright
\bibitem{N3} Max 2-SAT and the domination rule: results and ideas so far.
 Working note, 5 October 2026 (\texttt{max2sat\_results\_and\_ideas.tex}).
\bibitem{N4} Max 2-SAT with pruning clauses: pair decisions, sums, and a
 measure. Working note, 6--7 October 2026
 (\texttt{max2sat\_pairs\_sums\_measure.tex}).
\bibitem{N7} A simplified plan for Max 2-SAT: one search, one count, one Key
 Lemma. Working note, 8 October 2026 (\texttt{max2sat\_plan\_key\_lemma.tex}).
\bibitem{N8} The Key Lemma for random signs. Working note, 8 October 2026
 (\texttt{max2sat\_key\_lemma\_proof.tex}).
\bibitem{BvH16} A.~S.~Bandeira and R.~van Handel. Sharp nonasymptotic bounds
 on the norm of random matrices with independent entries. \emph{The Annals of
 Probability} 44(4):2479--2506, 2016.
\bibitem{Knu75} D.~E.~Knuth. Estimating the efficiency of backtrack
 programs. \emph{Mathematics of Computation} 29(129):121--136, 1975.
\bibitem{Mau03} A.~Maurer. A bound on the deviation probability for sums of
 non-negative random variables. \emph{Journal of Inequalities in Pure and
 Applied Mathematics} 4(1), Article~15, 2003.
\bibitem{Sza76} S.~J.~Szarek. On the best constants in the Khinchin inequality.
 \emph{Studia Mathematica} 58(2):197--208, 1976.
\bibitem{Tal95} M.~Talagrand. Concentration of measure and isoperimetric
 inequalities in product spaces. \emph{Publications Math\'ematiques de
 l'IH\'ES} 81:73--205, 1995.
\end{thebibliography}

\end{document}
